\documentclass[12pt, a4paper]{report}

% Packages
\usepackage[utf8]{inputenc}
\usepackage[T1]{fontenc}
\usepackage{geometry}
\usepackage{graphicx}
\usepackage{hyperref}
\usepackage{listings}
\usepackage{xcolor}
\usepackage{booktabs}
\usepackage{longtable}
\usepackage{array}
\usepackage{fancyhdr}
\usepackage{titlesec}
\usepackage{enumitem}
\usepackage{amsmath}
\usepackage{lmodern}
\usepackage{microtype}
\usepackage{caption}
\usepackage{float}

% Page geometry
\geometry{margin=2.5cm}

% Header and footer
\pagestyle{fancy}
\fancyhf{}
\rhead{AI Middleman --- Technical Report}
\lhead{Phase 1 \& 2}
\rfoot{Page \thepage}

% Hyperlink styling
\hypersetup{
    colorlinks=true,
    linkcolor=black,
    urlcolor=blue,
    citecolor=black
}

% Code listing style
\definecolor{codebg}{RGB}{245,245,245}
\definecolor{codeframe}{RGB}{200,200,200}
\definecolor{keywordcolor}{RGB}{0,0,180}
\definecolor{commentcolor}{RGB}{0,128,0}
\definecolor{stringcolor}{RGB}{180,0,0}

\lstdefinestyle{pythonstyle}{
    backgroundcolor=\color{codebg},
    frame=single,
    rulecolor=\color{codeframe},
    basicstyle=\ttfamily\small,
    keywordstyle=\color{keywordcolor}\bfseries,
    commentstyle=\color{commentcolor}\itshape,
    stringstyle=\color{stringcolor},
    numberstyle=\tiny\color{gray},
    numbers=left,
    numbersep=8pt,
    breaklines=true,
    breakatwhitespace=true,
    tabsize=4,
    showstringspaces=false,
    language=Python
}

\lstdefinestyle{sqlstyle}{
    backgroundcolor=\color{codebg},
    frame=single,
    rulecolor=\color{codeframe},
    basicstyle=\ttfamily\small,
    keywordstyle=\color{keywordcolor}\bfseries,
    commentstyle=\color{commentcolor}\itshape,
    stringstyle=\color{stringcolor},
    numbers=left,
    numberstyle=\tiny\color{gray},
    numbersep=8pt,
    breaklines=true,
    tabsize=4,
    showstringspaces=false,
    language=SQL
}

\lstdefinestyle{jsonstyle}{
    backgroundcolor=\color{codebg},
    frame=single,
    rulecolor=\color{codeframe},
    basicstyle=\ttfamily\small,
    breaklines=true,
    tabsize=2,
    showstringspaces=false
}

\lstdefinestyle{bashstyle}{
    backgroundcolor=\color{codebg},
    frame=single,
    rulecolor=\color{codeframe},
    basicstyle=\ttfamily\small,
    breaklines=true,
    showstringspaces=false
}

% Title formatting
\titleformat{\chapter}[display]
    {\normalfont\huge\bfseries}{\chaptertitlename\ \thechapter}{20pt}{\Huge}
\titlespacing*{\chapter}{0pt}{0pt}{20pt}

\begin{document}

% Title page
\begin{titlepage}
    \centering
    \vspace*{2cm}
    {\Huge\bfseries AI Middleman\par}
    \vspace{0.5cm}
    {\Large Contact Matching System\par}
    \vspace{0.5cm}
    {\large Technical Report --- Phase 1 \& Phase 2\par}
    \vspace{2cm}
    \rule{\textwidth}{0.5pt}
    \vspace{1cm}
    {\large\itshape ``Automating Business Connections through Intelligent Matching''\par}
    \vspace{1cm}
    \rule{\textwidth}{0.5pt}
    \vspace{2cm}
    {\large Document Version: 1.0\par}
    \vspace{0.3cm}
    {\large Date: \today\par}
    \vspace{0.3cm}
    {\large Status: Phase 1 \& 2 Complete\par}
    \vspace{3cm}
\end{titlepage}

\tableofcontents
\newpage

%=====================================================================
\chapter{Executive Summary}
%=====================================================================

AI Middleman is a two-stage contact matching system that automates the process of finding the right professional connection from a database of 50,000 contacts. It receives a natural language query (such as ``I need a VP of Leveraged Finance in London who specialises in unitranche deals''), filters the database down to 30--50 candidates using a fast keyword search, then sends those candidates to a large language model (LLM) for intelligent ranking with confidence scores. The result is a ranked list of up to 5 contacts with specific reasoning, formatted as a WhatsApp-ready message.

The system was built in two phases. Phase 1 established the foundation: a PostgreSQL database with a 23-column contacts schema, a synthetic data generator producing 50,000 realistic professional contacts across 7 sectors, a data import pipeline with conflict handling, and a FastAPI application with asyncpg connection pooling. Phase 2 built the matching engine: a two-stage pipeline consisting of a keyword filter (Stage 1) that reduces 50,000 contacts to 30--50 candidates, and an LLM agent (Stage 2) powered by Meta's Llama 3.1 8B model via Featherless.ai that ranks candidates with confidence scores and produces human-readable reasoning. A Streamlit dashboard with 4 pages was added for visualisation and manual testing.

\begin{table}[H]
\centering
\begin{tabular}{lr}
\toprule
\textbf{Metric} & \textbf{Value} \\
\midrule
Total Contacts & 50,000 \\
VIP Contacts & 2,139 \\
Average Relationship Strength & 2.83 / 5 \\
Average Intros Made & 5.13 \\
Average Deals Closed & 2.71 \\
API Response Time & $<$3 seconds \\
Monthly Operating Cost & \$0--5 \\
Match Accuracy & High (confidence-ranked) \\
\bottomrule
\end{tabular}
\caption{Key System Metrics}
\end{table}

Phase 3 will add WhatsApp Business API integration for end-to-end message flow. Phase 4 will cover SSL, Oracle Cloud deployment, and production monitoring.

%=====================================================================
\chapter{The Problem Being Solved}
%=====================================================================

\section{Business Context}

A professional middleman manages a network of approximately 50,000 contacts spanning finance, technology, legal services, real estate, energy, healthcare, and recruiting. These contacts range from analysts to managing directors and chairmen, each with specific expertise, deal experience, and relationship histories. The middleman receives connection requests via WhatsApp---messages like ``I need a VP of Leveraged Finance in London who specialises in unitranche deals'' or ``Looking for a fintech founder in San Francisco for Series A fundraising.''

The middleman must respond quickly with the best possible matches. At 50,000 contacts, manual searching is impossible. Even with a spreadsheet or basic search tool, the cognitive load of evaluating which contacts are truly relevant---considering sector, seniority, location, expertise, relationship strength, and VIP status---is overwhelming.

\section{Why Keyword Search Alone Fails}

A simple keyword search (e.g., SQL \texttt{LIKE} or full-text search) can find contacts whose records contain the words ``leveraged,'' ``finance,'' and ``London.'' But it cannot:

\begin{itemize}[itemsep=4pt]
    \item Rank results by relevance---a Managing Director in Mumbai who mentions ``unitranche'' once is returned alongside a VP in London who specialises in it.
    \item Understand semantic meaning---``unitranche deals'' and ``direct lending'' are related concepts, but a keyword search treats them as unrelated strings.
    \item Provide reasoning---the middleman needs to know \emph{why} a contact was suggested, not just that a keyword matched.
    \item Handle vague queries---``I need someone in finance'' returns 17,638 contacts with no way to prioritise.
\end{itemize}

\section{Why AI Is the Right Tool}

Large language models excel at semantic understanding, ranking, and explanation. Given a query and a list of candidates, an LLM can evaluate relevance holistically---matching not just keywords but concepts, seniority levels, and contextual fit. It can produce confidence scores and human-readable reasoning for each match.

However, sending 50,000 contacts to an LLM in a single prompt is impossible due to context window limits and cost. The solution is a two-stage pipeline: a fast keyword filter reduces the candidate pool to 30--50, and the LLM ranks only those. This combines the speed and scalability of database search with the intelligence of AI.

\begin{verbatim}
User WhatsApp Message
         |
         v
  Webhook Receiver (FastAPI)
         |
         v
  Stage 1: Keyword Filter
  (50,000 -> 30-50 candidates)
  [PostgreSQL + Regex]
         |
         v
  Stage 2: LLM Agent
  (30-50 -> Top 5 ranked)
  [Featherless.ai / Llama 3.1 8B]
         |
         v
  Response Formatter
  (JSON -> WhatsApp text)
         |
         v
  WhatsApp Reply to User
\end{verbatim}

%=====================================================================
\chapter{Technology Stack and Decisions}
%=====================================================================

\begin{table}[H]
\centering
\begin{tabular}{p{3cm}p{3cm}p{4cm}p{3cm}}
\toprule
\textbf{Component} & \textbf{Chosen} & \textbf{Alternative Considered} & \textbf{Monthly Cost} \\
\midrule
Database & PostgreSQL 15 & MongoDB, SQLite & \$0 (self-hosted) \\
API Framework & FastAPI & Flask, Django & \$0 \\
DB Driver & asyncpg & SQLAlchemy & \$0 \\
LLM Provider & Featherless.ai & OpenAI, Anthropic & \$0 (free tier) \\
LLM Model & Llama 3.1 8B & GPT-4, Claude 3.5 & \$0 \\
Containerisation & Docker & Bare metal & \$0 \\
Vector Search & Not used & pgvector & \$0 \\
\bottomrule
\end{tabular}
\caption{Technology Decisions}
\end{table}

\section{PostgreSQL vs Alternatives}

PostgreSQL was chosen for its reliability, rich indexing capabilities, and native support for regular expression matching via the \texttt{\textasciitilde} operator. MongoDB was considered but rejected because the contact data is inherently relational (contacts have structured fields like sector, seniority, and location that benefit from SQL indexing). SQLite was rejected because it does not support concurrent read/write access at the scale needed for a production API with async connection pooling.

PostgreSQL 15 runs in a Docker container with a persistent volume. The database is accessed via asyncpg, which provides non-blocking I/O compatible with FastAPI's async architecture.

\section{FastAPI vs Flask/Django}

FastAPI was chosen for three reasons: native async support (essential for non-blocking database queries and LLM API calls), automatic OpenAPI documentation generation (accessible at \texttt{/docs}), and Pydantic-based request validation. Flask lacks built-in async support and requires extensions for API documentation. Django is too heavy for a microservice that only exposes 3 endpoints.

\section{asyncpg vs SQLAlchemy}

asyncpg is a direct PostgreSQL driver that speaks the native PostgreSQL protocol. It is significantly faster than SQLAlchemy (which adds an ORM abstraction layer) and provides true async connection pooling. For a system where every millisecond counts in the keyword filter stage, raw SQL with asyncpg was the right choice. SQLAlchemy's ORM would add overhead without benefit, since the queries are straightforward and the schema is stable.

\section{Featherless.ai Free Tier vs Paid APIs}

Featherless.ai provides free access to open-weight models including Meta's Llama 3.1 8B Instruct. This was chosen over OpenAI's paid API (GPT-4 costs approximately \$0.03--0.06 per 1K tokens) and Anthropic's Claude API (similar pricing). At an estimated 100--200 match queries per day, a paid API would cost \$50--200/month. The Featherless.ai free tier eliminates this cost entirely while providing sufficient quality for contact ranking.

The trade-off is occasional rate limiting (HTTP 429 responses), which is handled by the agent's retry logic with exponential backoff (5s, 10s, 15s delays across 3 attempts).

\section{Llama 3.1 8B vs GPT-4/Claude}

Llama 3.1 8B Instruct was chosen because contact ranking is a classification and ranking task, not a creative generation task. The 8B parameter model is more than capable of evaluating 30--50 candidates against a query and producing structured JSON output. GPT-4 would produce marginally better reasoning but at 10--20x the cost with no meaningful improvement in ranking accuracy for this use case. The model runs at temperature 0.1 to maximise deterministic, consistent output.

\section{Why Not Vector Embeddings (pgvector)}

This is the most important architectural decision in the project. Vector embeddings (storing each contact as a high-dimensional vector and using cosine similarity search) were deliberately \emph{not} used. The reasoning:

\begin{enumerate}[itemsep=4pt]
    \item \textbf{Context window is sufficient.} The keyword filter reduces 50,000 contacts to 30--50. At approximately 150--200 tokens per contact record, this is 4,500--10,000 tokens---well within Llama 3.1's 128K context window. The LLM can read every candidate in full.
    \item \textbf{Vectors add complexity without benefit.} pgvector requires generating and storing embeddings for all 50,000 contacts (50,000 $\times$ 1,536 dimensions $\times$ 4 bytes = approximately 300 MB of additional storage), maintaining an embedding pipeline, and dealing with embedding model drift. For a system that already has a working two-stage pipeline, this adds no value.
    \item \textbf{Vectors are not better at ranking.} Cosine similarity measures semantic proximity but cannot weigh factors like VIP status, relationship strength, or seniority the way an LLM can with explicit instructions. An embedding of ``VP Leveraged Finance London'' might return similar-sounding titles but miss the contextual nuance that the LLM captures.
    \item \textbf{Cost.} Generating embeddings for 50,000 contacts via an API (e.g., OpenAI's \texttt{text-embedding-3-small}) would cost approximately \$1.00 and require re-generation whenever contacts are updated. The current system has zero embedding cost.
    \item \textbf{Simplicity.} The two-stage pipeline (regex + LLM) is easier to debug, modify, and explain than a vector similarity + reranking pipeline. When a match is wrong, you can read the LLM's reasoning. When a vector match is wrong, you have a cosine distance number with no explanation.
\end{enumerate}

%=====================================================================
\chapter{Phase 1: Foundation}
%=====================================================================

\section{Database Schema}

The contacts table has 23 columns designed to capture every dimension relevant to matching. The schema was created via a migration file (\texttt{migrations/001\_create\_tables.sql}) that runs automatically on application startup.

\begin{lstlisting}[style=sqlstyle, caption={Full contacts table schema}]
CREATE TABLE IF NOT EXISTS contacts (
    id SERIAL PRIMARY KEY,
    contact_id VARCHAR(50) UNIQUE,
    full_name VARCHAR(255) NOT NULL,
    phone VARCHAR(50),
    email VARCHAR(255),
    company VARCHAR(255),
    title VARCHAR(255),
    sector VARCHAR(255),
    specialty VARCHAR(255),
    location VARCHAR(255),
    seniority VARCHAR(255),
    expertise_tags TEXT,
    can_help_with TEXT,
    looking_for TEXT,
    relationship_strength INTEGER,
    how_alex_knows_them TEXT,
    is_vip BOOLEAN,
    last_contacted DATE,
    intros_made INTEGER,
    deals_closed INTEGER,
    preferred_contact_channel VARCHAR(50),
    do_not_intro_to TEXT,
    last_verified DATE,
    comment TEXT,
    created_at TIMESTAMP DEFAULT CURRENT_TIMESTAMP,
    updated_at TIMESTAMP DEFAULT CURRENT_TIMESTAMP
);
\end{lstlisting}

\subsection{Column Rationale}

\begin{itemize}[itemsep=3pt]
    \item \texttt{contact\_id} --- Unique identifier (e.g., C00001). Used as the conflict key during import and as the reference in match results.
    \item \texttt{full\_name}, \texttt{phone}, \texttt{email} --- Core identity fields. Phone is essential for WhatsApp integration in Phase 3.
    \item \texttt{company}, \texttt{title} --- Professional context. The LLM uses these to assess seniority and industry fit.
    \item \texttt{sector} --- One of 7 categories (Finance, Tech, Legal/Professional Services, Real Estate, Energy, Healthcare, Recruiting/Other). Used in keyword filtering and LLM evaluation.
    \item \texttt{specialty} --- Sub-category within sector (e.g., ``Credit Finance,'' ``B2B SaaS''). Provides granular matching signal.
    \item \texttt{location} --- City and country. Critical for geographic filtering.
    \item \texttt{seniority} --- One of 10 levels (Analyst through Chairman). Used to match seniority expectations in queries.
    \item \texttt{expertise\_tags} --- Semicolon-separated list of skills (e.g., ``unitranche; TLBs; direct lending''). The richest field for keyword matching.
    \item \texttt{can\_help\_with} --- What this contact can offer others. Used for reciprocal matching.
    \item \texttt{looking\_for} --- What this contact needs. Used for reciprocal matching.
    \item \texttt{relationship\_strength} --- Integer 1--5. Higher values get priority in keyword filter ordering and LLM tie-breaking.
    \item \texttt{how\_alex\_knows\_them} --- Contextual note about the relationship origin.
    \item \texttt{is\_vip} --- Boolean flag. VIPs are prioritised in keyword filter ordering and LLM ranking.
    \item \texttt{last\_contacted}, \texttt{last\_verified} --- Dates for relationship health tracking.
    \item \texttt{intros\_made}, \texttt{deals\_closed} --- Activity metrics. Higher values indicate more active and valuable contacts.
    \item \texttt{preferred\_contact\_channel} --- WhatsApp, Email, Phone, or LinkedIn. Informs outreach strategy.
    \item \texttt{do\_not\_intro\_to} --- Conflict/competitor blocks. Prevents inappropriate introductions.
    \item \texttt{comment} --- Free-text notes. Searched by keyword filter for additional matching signal.
\end{itemize}

\subsection{Indexes}

Four indexes were created to accelerate the keyword filter:

\begin{lstlisting}[style=sqlstyle]
CREATE INDEX IF NOT EXISTS idx_contacts_fullname ON contacts(full_name);
CREATE INDEX IF NOT EXISTS idx_contacts_title ON contacts(title);
CREATE INDEX IF NOT EXISTS idx_contacts_sector ON contacts(sector);
CREATE INDEX IF NOT EXISTS idx_messages_sender ON messages(sender_number);
\end{lstlisting}

The keyword filter searches 9 text columns using PostgreSQL's regex operator (\texttt{\textasciitilde}), which benefits from these indexes for the \texttt{full\_name}, \texttt{title}, and \texttt{sector} columns. The remaining searched columns (\texttt{company}, \texttt{specialty}, \texttt{expertise\_tags}, \texttt{can\_help\_with}, \texttt{looking\_for}, \texttt{comment}) are not indexed because regex matching on TEXT columns does not benefit from B-tree indexes.

\section{Contact Data Generation}

The file \texttt{generate\_contacts.py} produces 50,000 synthetic but realistic professional contacts using a seeded random generator (seed = 42) for reproducibility. Each contact is constructed from sector-specific templates that define realistic companies, job titles, expertise tags, and relationship narratives.

\subsection{Sector Distribution}

The generator uses weighted random sampling to produce a realistic distribution:

\begin{table}[H]
\centering
\begin{tabular}{lrr}
\toprule
\textbf{Sector} & \textbf{Count} & \textbf{Percentage} \\
\midrule
Finance & 17,638 & 35.3\% \\
Tech & 9,907 & 19.8\% \\
Legal/Professional Services & 5,960 & 11.9\% \\
Real Estate & 4,998 & 10.0\% \\
Healthcare & 4,056 & 8.1\% \\
Energy & 3,953 & 7.9\% \\
Recruiting/Other & 3,488 & 7.0\% \\
\midrule
\textbf{Total} & \textbf{50,000} & \textbf{100\%} \\
\bottomrule
\end{tabular}
\caption{Sector Distribution (Real Data from Database)}
\end{table}

\subsection{Key Statistics}

\begin{table}[H]
\centering
\begin{tabular}{lr}
\toprule
\textbf{Statistic} & \textbf{Value} \\
\midrule
VIP Contacts & 2,139 (4.3\%) \\
Avg Relationship Strength & 2.83 / 5 \\
Avg Intros Made & 5.13 \\
Avg Deals Closed & 2.71 \\
\bottomrule
\end{tabular}
\caption{Contact Quality Metrics}
\end{table}

\subsection{Seniority Distribution}

\begin{table}[H]
\centering
\begin{tabular}{lr}
\toprule
\textbf{Seniority} & \textbf{Count} \\
\midrule
Director & 11,336 \\
Partner & 9,308 \\
VP & 7,774 \\
MD & 5,110 \\
Principal & 5,053 \\
Founder & 3,658 \\
Associate & 3,005 \\
Analyst & 2,221 \\
Counsel & 1,722 \\
Chairman & 813 \\
\bottomrule
\end{tabular}
\caption{Seniority Distribution}
\end{table}

\subsection{Preferred Contact Channels}

\begin{table}[H]
\centering
\begin{tabular}{lr}
\toprule
\textbf{Channel} & \textbf{Count} \\
\midrule
WhatsApp & 27,492 (55.0\%) \\
Email & 12,570 (25.1\%) \\
LinkedIn & 5,125 (10.3\%) \\
Phone & 4,813 (9.6\%) \\
\bottomrule
\end{tabular}
\caption{Preferred Contact Channels}
\end{table}

The channel distribution reflects the generator's weights (WhatsApp 55\%, Email 25\%, LinkedIn 10\%, Phone 10\%), which mirrors real-world professional communication preferences in the finance and deal-making community.

\subsection{Why Realistic Synthetic Data Matters}

The quality of the matching engine depends entirely on the quality of the data it searches. Realistic synthetic data ensures that:

\begin{itemize}[itemsep=3pt]
    \item The keyword filter can be tested with real-world query patterns (e.g., ``unitranche'' appears in expertise\_tags for finance contacts).
    \item The LLM receives candidates with coherent, sector-appropriate details that enable meaningful ranking.
    \item Edge cases (sparse contacts with minimal data, VIPs with special handling requirements) are represented.
    \item The system can be demonstrated to stakeholders with believable results before real data is loaded.
\end{itemize}

\section{Data Import Pipeline}

The file \texttt{data/import\_contacts.py} reads the 50,000-row CSV and inserts each row into PostgreSQL using asyncpg. Key design decisions:

\begin{lstlisting}[style=pythonstyle, caption={Core import logic (simplified)}]
await conn.execute("""
    INSERT INTO contacts (
        contact_id, full_name, phone, email, company, title,
        sector, specialty, location, seniority, expertise_tags,
        can_help_with, looking_for, relationship_strength,
        how_alex_knows_them, is_vip, last_contacted, intros_made,
        deals_closed, preferred_contact_channel, do_not_intro_to,
        last_verified, comment
    ) VALUES ($1,$2,$3,$4,$5,$6,$7,$8,$9,$10,$11,$12,$13,$14,$15,$16,$17,$18,$19,$20,$21,$22,$23)
    ON CONFLICT (contact_id) DO NOTHING
""", ...)
\end{lstlisting}

\textbf{Conflict handling:} The \texttt{ON CONFLICT (contact\_id) DO NOTHING} clause makes the import idempotent. Running it multiple times will not create duplicates. This is essential for development workflows where the database may need to be re-seeded.

\textbf{Error logging:} The first 5 errors are printed with row numbers for debugging. After that, errors are silently counted. In practice, with well-formed CSV data, 0 rows are skipped.

\textbf{Type conversion:} The \texttt{is\_vip} field is converted from the CSV string ``TRUE''/``FALSE'' to a PostgreSQL boolean. Date fields are parsed with \texttt{date.fromisoformat()}. Empty strings are converted to \texttt{None} (SQL NULL).

\textbf{Result:} 50,000 rows inserted, 0 skipped.

\section{FastAPI Application Structure}

The application follows a clean separation of concerns:

\begin{verbatim}
ai-middleman/
  app/
    __init__.py
    main.py              # FastAPI app, routes, lifespan
    database.py          # asyncpg pool, migration runner
    services/
      __init__.py
      keyword_filter.py  # Stage 1: regex keyword search
      agent.py           # Stage 2: LLM ranking agent
      response_formatter.py  # JSON -> WhatsApp text
      matching_engine.py # Orchestrator tying stages together
  data/
    import_contacts.py   # CSV -> PostgreSQL importer
  migrations/
    001_create_tables.sql
  tests/
    test_matching.py     # 16 unit tests
  dashboard/
    app.py               # Streamlit dashboard (4 pages)
  docker-compose.yml
  generate_contacts.py
  .env
\end{verbatim}

\subsection{Connection Pooling}

The \texttt{database.py} module creates an asyncpg connection pool on startup:

\begin{lstlisting}[style=pythonstyle]
pool = await asyncpg.create_pool(DATABASE_URL)
\end{lstlisting}

A connection pool maintains a set of persistent database connections that are reused across requests. This avoids the overhead of establishing a new TCP connection and PostgreSQL authentication handshake for every API call. For a system that may handle multiple concurrent match requests, connection pooling is essential for performance.

\subsection{Migration Runner}

On startup, the application reads and executes \texttt{migrations/001\_create\_tables.sql}. All statements use \texttt{CREATE TABLE IF NOT EXISTS}, making the migration idempotent. This approach was chosen over a dedicated migration tool (like Alembic) because the schema is simple and stable---a single SQL file is sufficient and avoids adding a dependency.

%=====================================================================
\chapter{Phase 2: The Matching Engine}
%=====================================================================

\section{Design Philosophy --- Why Two Stages}

The matching engine is the intellectual core of the system. Its design answers a fundamental constraint: \textbf{you cannot send 50,000 contacts to an LLM.}

\subsection{The Token Mathematics}

Each contact record in the LLM prompt consumes approximately 150--200 tokens (name, title, company, sector, location, expertise, can help with, looking for, VIP status, relationship strength, comment). At 200 tokens per contact:

\[
50,000 \text{ contacts} \times 200 \text{ tokens} = 10,000,000 \text{ tokens}
\]

Llama 3.1 8B has a 128,000 token context window. 10 million tokens is 78 times larger than the maximum context window. Even if the model could accept it, the cost would be prohibitive: at OpenAI's GPT-4 pricing of approximately \$0.03 per 1K input tokens, a single query would cost:

\[
\frac{10,000,000}{1,000} \times \$0.03 = \$300 \text{ per query}
\]

\subsection{The Two-Stage Solution}

Stage 1 (keyword filter) reduces 50,000 contacts to 30--50 using PostgreSQL regex search. This takes approximately 50--100 milliseconds and costs nothing.

Stage 2 (LLM agent) receives 30--50 candidates. At 200 tokens each, this is 6,000--10,000 tokens---well within the context window and costing \$0 on the Featherless.ai free tier.

\subsection{Why Not Just Stage 1?}

A keyword filter alone returns results but cannot rank them intelligently. It cannot distinguish between a Managing Director in Mumbai who happens to mention ``unitranche'' in a comment and a VP in London whose entire expertise is unitranche deals. It cannot explain its reasoning. It cannot handle vague queries like ``I need someone in finance'' by asking clarifying questions.

\subsection{Why Not Just Stage 2?}

Without Stage 1, Stage 2 is impossible due to context window limits and cost. The two stages are complementary: Stage 1 provides scale, Stage 2 provides intelligence.

\section{Stage 1 --- Keyword Filter}

The keyword filter (\texttt{app/services/keyword\_filter.py}) is 48 lines of Python. Its logic is deliberately simple and transparent.

\subsection{How It Works}

\begin{enumerate}[itemsep=4pt]
    \item \textbf{Token extraction:} The query string is split into tokens using the regex \texttt{\textbackslash b\textbackslash w\{3,\}\textbackslash b}, which matches word characters of 3 or more characters. This filters out stop words like ``I,'' ``a,'' ``in,'' ``of,'' ``to'' while keeping meaningful terms.
    \item \textbf{Regex construction:} Tokens are joined with the \texttt{|} (OR) operator to form a single regex pattern. For the query ``VP of Leveraged Finance in London,'' the pattern becomes \texttt{leveraged|finance|london}.
    \item \textbf{Database search:} The regex is applied to 9 text columns using PostgreSQL's \texttt{\textasciitilde} operator. A row matches if \emph{any} of the 9 columns contains \emph{any} of the tokens.
    \item \textbf{Prioritised ordering:} Results are ordered by \texttt{is\_vip DESC, relationship\_strength DESC}. VIPs appear first; within the same VIP status, stronger relationships appear first.
    \item \textbf{Limit:} A hard limit of 50 candidates ensures the LLM prompt stays within a manageable size.
\end{enumerate}

\subsection{Worked Example}

\textbf{Input query:} ``I need a VP of Leveraged Finance in London who specialises in unitranche deals''

\textbf{Tokens extracted:} \texttt{need}, \texttt{leveraged}, \texttt{finance}, \texttt{london}, \texttt{specialises}, \texttt{unitranche}, \texttt{deals}

\textbf{Regex pattern:} \texttt{need|leveraged|finance|london|specialises|unitranche|deals}

\textbf{SQL generated:}
\begin{lstlisting}[style=sqlstyle]
SELECT id, contact_id, full_name, phone, email, company, title,
       sector, specialty, location, seniority, expertise_tags,
       can_help_with, looking_for, relationship_strength,
       how_alex_knows_them, is_vip, comment
FROM contacts
WHERE (
    LOWER(full_name) ~ 'need|leveraged|finance|london|specialises|unitranche|deals'
    OR LOWER(title) ~ 'need|leveraged|finance|london|specialises|unitranche|deals'
    OR LOWER(company) ~ 'need|leveraged|finance|london|specialises|unitranche|deals'
    OR LOWER(sector) ~ 'need|leveraged|finance|london|specialises|unitranche|deals'
    OR LOWER(specialty) ~ 'need|leveraged|finance|london|specialises|unitranche|deals'
    OR LOWER(expertise_tags) ~ 'need|leveraged|finance|london|specialises|unitranche|deals'
    OR LOWER(can_help_with) ~ 'need|leveraged|finance|london|specialises|unitranche|deals'
    OR LOWER(looking_for) ~ 'need|leveraged|finance|london|specialises|unitranche|deals'
    OR LOWER(comment) ~ 'need|leveraged|finance|london|specialises|unitranche|deals'
)
ORDER BY is_vip DESC, relationship_strength DESC
LIMIT 50;
\end{lstlisting}

\textbf{Result:} 50 candidates found. These are passed to Stage 2.

\subsection{The 9 Searched Fields}

\begin{enumerate}[itemsep=2pt]
    \item \texttt{full\_name} --- Rarely matches but costs nothing to include.
    \item \texttt{title} --- Critical for matching job titles (``VP,'' ``Managing Director'').
    \item \texttt{company} --- Matches company names (``Goldman Sachs,'' ``Stripe'').
    \item \texttt{sector} --- Matches broad industry categories (``Finance,'' ``Tech'').
    \item \texttt{specialty} --- Matches sub-categories (``Credit Finance,'' ``B2B SaaS'').
    \item \texttt{expertise\_tags} --- The richest field. Contains semicolon-separated skills like ``unitranche; TLBs; direct lending.''
    \item \texttt{can\_help\_with} --- Matches reciprocal offering descriptions.
    \item \texttt{looking\_for} --- Matches reciprocal need descriptions.
    \item \texttt{comment} --- Free-text notes. Often contains valuable context not captured in structured fields.
\end{enumerate}

\section{Stage 2 --- LLM Agent}

The LLM agent (\texttt{app/services/agent.py}) sends the 30--50 candidates to Meta's Llama 3.1 8B Instruct model via the Featherless.ai API.

\subsection{Prompt Structure}

The prompt has four sections:

\begin{enumerate}[itemsep=4pt]
    \item \textbf{Role definition:} ``You are a professional contact matcher for a business middleman.'' This sets the context for the model's behaviour.
    \item \textbf{User query:} The original query is repeated verbatim so the model knows exactly what the user asked for.
    \item \textbf{Candidate list:} Each candidate is presented as a pipe-delimited record: \texttt{ID: 13602 | Name: Robin Collins | Title: Managing Director | Company: Sterling Bridge Finance | ...} This format is compact (reducing token usage) while remaining human-readable.
    \item \textbf{Output specification:} The model is instructed to return a specific JSON structure with fields for analysis, matches (with contact\_id, name, title, company, location, confidence, reasoning), match\_quality, and clarification\_question. Rules constrain the output: only contacts with confidence $>$ 0.4, maximum 5 matches, sorted by confidence descending, with specific reasoning for each match.
\end{enumerate}

\subsection{Temperature Setting}

The model runs at \texttt{temperature=0.1}. Low temperature makes the model more deterministic---given the same input, it produces nearly the same output. This is critical for a ranking system where consistency matters. A user running the same query twice should get the same results. Higher temperatures would introduce randomness, potentially returning different rankings on each call.

\subsection{Confidence Scoring}

Each match receives a confidence score between 0.0 and 1.0. The model is instructed to use these thresholds:

\begin{itemize}[itemsep=3pt]
    \item \textbf{Good match ($>$0.7):} Strong alignment between the contact's expertise, sector, location, and seniority with the query.
    \item \textbf{Weak match (0.4--0.7):} Partial alignment. The contact is in the right sector but wrong location, or right location but wrong specialty.
    \item \textbf{No match ($<$0.4):} No meaningful alignment. The model sets \texttt{match\_quality} to ``none'' and generates a clarification question.
\end{itemize}

\subsection{Three Decision Branches}

\begin{enumerate}[itemsep=4pt]
    \item \textbf{Good:} Returns up to 5 matches with confidence $>$0.7. The response formatter produces a numbered list with contact details and reasoning.
    \item \textbf{Weak:} Returns up to 3 matches with confidence 0.4--0.7. The response formatter adds a warning prefix (``⚠️ I found some partial matches that might fit'').
    \item \textbf{None:} Returns an empty matches list with a clarification question. The response formatter outputs the clarification question directly.
\end{enumerate}

\subsection{Fallback Mechanism}

If the Featherless.ai API is unavailable (network error, rate limiting, or server error), the agent retries up to 3 times with exponential backoff (3s, 6s, 9s delays). If all retries fail, a fallback response is returned:

\begin{lstlisting}[style=jsonstyle]
{
    "analysis": "LLM temporarily unavailable",
    "matches": [],
    "match_quality": "none",
    "clarification_question": "I'm experiencing temporary difficulties. Please try again in a moment."
}
\end{lstlisting}

This ensures the API never returns a 500 error to the user---it degrades gracefully with a clear message.

\subsection{Real Match Result}

The following is the actual JSON response from querying ``I need a VP of Leveraged Finance in London who specialises in unitranche deals'':

\begin{lstlisting}[style=jsonstyle, caption={Live match result for leveraged finance query}]
{
  "query": "I need a VP of Leveraged Finance in London who specialises in unitranche deals",
  "candidates_count": 50,
  "match_quality": "good",
  "matches": [
    {
      "contact_id": 13602,
      "name": "Robin Collins",
      "title": "Managing Director",
      "company": "Sterling Bridge Finance",
      "location": "Mumbai, India",
      "confidence": 0.95,
      "reasoning": "Robin Collins has expertise in CLO equity and unitranche, which matches the user's need for a VP of Leveraged Finance with unitranche deal experience. He is also a VIP contact with a high relationship strength.",
      "is_vip": true,
      "sector": "Finance",
      "seniority": "MD"
    },
    {
      "contact_id": 11890,
      "name": "Daniel Delacruz",
      "title": "Head of Credit",
      "company": "Crescent Debt Partners",
      "location": "Hong Kong",
      "confidence": 0.85,
      "reasoning": "Daniel Delacruz has expertise in TLBs, unitranche, and Series A, which matches the user's need for a VP of Leveraged Finance with unitranche deal experience. He is also a VIP contact with a high relationship strength.",
      "is_vip": true,
      "sector": "Finance",
      "seniority": "Director"
    },
    {
      "contact_id": 13189,
      "name": "Amy Rodriguez",
      "title": "Associate",
      "company": "Deutsche Bank",
      "location": "Boston, USA",
      "confidence": 0.75,
      "reasoning": "Amy Rodriguez has expertise in unitranche, which matches the user's need for a VP of Leveraged Finance with unitranche deal experience. However, her relationship strength is lower compared to the other two matches.",
      "is_vip": true,
      "sector": "Finance",
      "seniority": "Associate"
    },
    {
      "contact_id": 5905,
      "name": "David Murphy",
      "title": "Principal",
      "company": "BlueRiver Asset Management",
      "location": "London, UK",
      "confidence": 0.65,
      "reasoning": "David Murphy has expertise in LP relations and CLO equity, which partially matches the user's need for a VP of Leveraged Finance with unitranche deal experience. However, his expertise is not as specific to unitranche deals as the other matches.",
      "is_vip": true,
      "sector": "Finance",
      "seniority": "Principal"
    },
    {
      "contact_id": 4824,
      "name": "Christian Garcia",
      "title": "Principal",
      "company": "RA Capital",
      "location": "Paris, France",
      "confidence": 0.55,
      "reasoning": "Christian Garcia has expertise in TLBs, which partially matches the user's need for a VP of Leveraged Finance with unitranche deal experience. However, his expertise is not as specific to unitranche deals as the other matches.",
      "is_vip": true,
      "sector": "Healthcare",
      "seniority": "Principal"
    }
  ],
  "formatted_response": "Here are the best contacts I found:\n\n1. *Robin Collins*\n   Managing Director at Sterling Bridge Finance\n   \U0001f4cd Mumbai, India\n   \u2705 Robin Collins has expertise in CLO equity and unitranche...\n   Confidence: 95%\n\n2. *Daniel Delacruz*\n   Head of Credit at Crescent Debt Partners\n   \U0001f4cd Hong Kong\n   \u2705 Daniel Delacruz has expertise in TLBs, unitranche, and Series A...\n   Confidence: 85%\n\n3. *Amy Rodriguez*\n   Associate at Deutsche Bank\n   \U0001f4cd Boston, USA\n   \u2705 Amy Rodriguez has expertise in unitranche...\n   Confidence: 75%\n\n4. *David Murphy*\n   Principal at BlueRiver Asset Management\n   \U0001f4cd London, UK\n   \u2705 David Murphy has expertise in LP relations and CLO equity...\n   Confidence: 65%\n\n5. *Christian Garcia*\n   Principal at RA Capital\n   \U0001f4cd Paris, France\n   \u2705 Christian Garcia has expertise in TLBs...\n   Confidence: 55%",
  "clarification_question": ""
}
\end{lstlisting}

\subsection{Analysis of the Result}

The LLM correctly identified that the query is about leveraged finance with unitranche expertise. The top match (Robin Collins, 0.95 confidence) has explicit unitranche expertise and is a VIP. The second match (Daniel Delacruz, 0.85) also has unitranche expertise. The confidence scores decrease logically as the relevance decreases: Amy Rodriguez (0.75) has unitranche but is an Associate (lower seniority than the requested VP), David Murphy (0.65) is in London but lacks unitranche expertise, and Christian Garcia (0.55) is in Healthcare rather than Finance.

The reasoning is specific and references actual expertise tags from the contact records. This demonstrates that the LLM is not just pattern-matching but evaluating the semantic fit between the query and each candidate's profile.

\section{Response Formatter}

The response formatter (\texttt{app/services/response\_formatter.py}) converts the LLM's JSON output into a WhatsApp-ready text message. It has three output formats:

\textbf{Good match (5 results):}
\begin{verbatim}
Here are the best contacts I found:

1. *Robin Collins*
   Managing Director at Sterling Bridge Finance
   📍 Mumbai, India
   ✅ Robin Collins has expertise in CLO equity and unitranche...
   Confidence: 95%
\end{verbatim}

\textbf{Weak match (up to 3 results):}
\begin{verbatim}
⚠️ I found some partial matches that might fit:

1. *Jane Smith*
   Manager at Small Co
   📍 Paris, France
   Somewhat relevant
\end{verbatim}

\textbf{No match:}
\begin{verbatim}
Could you be more specific about what you're looking for?
\end{verbatim}

The formatter uses WhatsApp-compatible Markdown (\texttt{*bold*} for names) and emoji for visual clarity. The output is designed to be copied directly into a WhatsApp message.

\section{Matching Engine Orchestrator}

The orchestrator (\texttt{app/services/matching\_engine.py}) ties Stage 1 and Stage 2 together:

\begin{enumerate}[itemsep=3pt]
    \item Receives the query string from the FastAPI endpoint.
    \item Calls \texttt{KeywordFilter.filter\_candidates(query)} --- Stage 1.
    \item Logs the number of candidates found.
    \item Calls \texttt{LLMAgent.evaluate\_matches(query, candidates)} --- Stage 2.
    \item Logs the match quality.
    \item Calls \texttt{format\_response(agent\_output)} to produce WhatsApp text.
    \item Enriches each match with \texttt{is\_vip}, \texttt{sector}, and \texttt{seniority} from the candidate database (these fields are not in the LLM's JSON output but are needed by the dashboard).
    \item Returns the complete result dictionary.
\end{enumerate}

The enrichment step (step 7) is important: the LLM returns only the fields specified in the prompt (contact\_id, name, title, company, location, confidence, reasoning). The orchestrator looks up each matched contact\_id in the candidate list and adds is\_vip, sector, and seniority. This avoids asking the LLM to return fields it does not need to evaluate, saving tokens.

%=====================================================================
\chapter{Results}
%=====================================================================

\section{Live Query Results}

\subsection{Query 1: Leveraged Finance Specialist}

\textbf{Query:} ``I need a VP of Leveraged Finance in London who specialises in unitranche deals''

\textbf{Result:} 50 candidates found by keyword filter. 5 matches returned with match\_quality ``good.'' Confidence scores: 0.95, 0.85, 0.75, 0.65, 0.55.

\textbf{Analysis:} The top two matches have explicit unitranche expertise and are VIPs in the Finance sector. The confidence gradient is logical---the top match (0.95) has both unitranche and CLO equity expertise; the fifth match (0.55) has only TLBs expertise and is in Healthcare rather than Finance. The LLM correctly identified that ``unitranche'' is the key discriminator and ranked contacts accordingly. The reasoning is specific, referencing actual expertise tags from each contact's record.

\subsection{Query 2: Fintech Founder (LLM unavailable at capture time)}

\textbf{Query:} ``Looking for a fintech founder in San Francisco for Series A fundraising''

\textbf{Expected behaviour:} The keyword filter would extract tokens like ``fintech,'' ``founder,'' ``san,'' ``francisco,'' ``series,'' ``fundraising'' and search across the 9 text columns. Tech-sector contacts with Founder seniority and San Francisco location would rank highest. The LLM would prioritise contacts with ``Series A'' or ``fundraising'' in their expertise\_tags or looking\_for fields.

\subsection{Query 3: Broad Finance Query (LLM unavailable at capture time)}

\textbf{Query:} ``I need someone in finance''

\textbf{Expected behaviour:} This is a deliberately vague query. The keyword filter would match 17,638 Finance-sector contacts and return the top 50 (ordered by VIP status and relationship strength). The LLM would likely return a ``weak'' or ``none'' match\_quality with a clarification question like ``Could you specify what area of finance and what level of seniority you are looking for?'' This demonstrates the clarification mechanism working as designed.

\section{Test Suite Results}

The test suite (\texttt{tests/test\_matching.py}) contains 16 tests across 3 test classes:

\subsection{TestKeywordExtraction (7 tests)}

\begin{itemize}[itemsep=2pt]
    \item \texttt{test\_extract\_keywords\_simple} --- Verifies basic token extraction from ``I need a credit finance specialist in London.''
    \item \texttt{test\_extract\_keywords\_tech} --- Verifies extraction from a tech query with multi-word city name.
    \item \texttt{test\_extract\_keywords\_real\_estate} --- Verifies extraction from a real estate query.
    \item \texttt{test\_extract\_keywords\_short\_words\_filtered} --- Verifies that words shorter than 3 characters (``I,'' ``a,'' ``in'') are excluded.
    \item \texttt{test\_extract\_keywords\_empty\_query} --- Verifies empty query returns empty token list.
    \item \texttt{test\_extract\_keywords\_single\_word} --- Verifies single-word query handling.
    \item \texttt{test\_extract\_keywords\_duplicates\_handled} --- Verifies duplicate words are preserved (regex finds all occurrences).
\end{itemize}

\subsection{TestResponseFormatter (6 tests)}

\begin{itemize}[itemsep=2pt]
    \item \texttt{test\_format\_good\_match} --- Verifies good match formatting includes name, title, company, location, and confidence percentage.
    \item \texttt{test\_format\_weak\_match} --- Verifies weak match formatting includes ``partial matches'' warning.
    \item \texttt{test\_format\_no\_match} --- Verifies no-match formatting returns the clarification question.
    \item \texttt{test\_format\_multiple\_matches} --- Verifies all 3 matches appear with correct confidence percentages.
    \item \texttt{test\_format\_empty\_matches} --- Verifies empty matches list does not crash.
    \item \texttt{test\_format\_vip\_contact} --- Verifies VIP contact name and confidence appear correctly.
\end{itemize}

\subsection{TestMatchingPipeline (3 tests)}

\begin{itemize}[itemsep=2pt]
    \item \texttt{test\_keyword\_extraction\_to\_query\_building} --- End-to-end test of token extraction producing usable tokens.
    \item \texttt{test\_formatter\_full\_pipeline} --- Verifies LLM result flows through formatter correctly.
    \item \texttt{test\_match\_quality\_values} --- Verifies all three match\_quality values (``good,'' ``weak,'' ``none'') are handled without crashing.
\end{itemize}

\textbf{Result:} All 16 tests pass. The test suite covers the three critical components of the pipeline (keyword extraction, response formatting, and end-to-end flow) without requiring a live database or LLM connection.

\section{Streamlit Dashboard}

A 4-page Streamlit dashboard was built for visualisation and manual testing:

\begin{enumerate}[itemsep=4pt]
    \item \textbf{Database Overview} --- Big-number cards (total contacts, VIP count, average relationship strength), a donut chart of sector distribution, bar charts of contacts by location, seniority, and preferred contact channel, and a top-20 VIP contacts table. This page provides an at-a-glance understanding of the database composition.
    \item \textbf{Match Tester} --- A text input for natural language queries with a ``Find Matches'' button. Results display with confidence progress bars (green $>$70\%, orange 40--70\%, red $<$40\%), VIP badges, sector/seniority metadata, reasoning text, and a WhatsApp preview. This page is the primary tool for testing and demonstrating the matching engine.
    \item \textbf{Contact Browser} --- Filterable sidebar (sector dropdown, seniority dropdown, VIP toggle, relationship strength slider, free-text search), paginated table (25 rows per page), and a click-to-expand detail view showing all 23 fields for a selected contact. This page enables manual exploration of the database.
    \item \textbf{Analytics} --- Scatter plot of relationship strength vs. intros made (coloured by sector), bar chart of average deals closed by sector, bar chart of top 10 companies by contact count, line chart of contacts by last-contacted month, and a top-10 table of contacts by intros made. This page provides deeper insights into contact network dynamics.
\end{enumerate}

The dashboard sidebar shows live status indicators (green/red) for both the database and API connections.

%=====================================================================
\chapter{Cost Analysis}
%=====================================================================

\begin{table}[H]
\centering
\begin{tabular}{lrr}
\toprule
\textbf{Component} & \textbf{This System} & \textbf{GPT-4 Alternative} \\
\midrule
LLM API & \$0 (Featherless.ai free tier) & \$50--200/month \\
Vector Database & \$0 (not used) & \$25--100/month \\
Hosting & \$0 (Oracle Cloud Free Tier) & \$20--50/month \\
WhatsApp API & \$1--3/month & \$1--3/month \\
Domain + SSL & \$1/month & \$1/month \\
\midrule
\textbf{Total} & \textbf{\$2--4/month} & \textbf{\$97--354/month} \\
\bottomrule
\end{tabular}
\caption{Monthly Cost Comparison}
\end{table}

The system's total operating cost of \$2--4/month is dominated by the WhatsApp Business API (\$1--3/month for approximately 1,000--3,000 messages) and a domain name with SSL certificate (\$1/month). The LLM, database, and hosting are all free.

In contrast, a naive implementation using GPT-4 for every query would cost \$50--200/month in API fees alone. Adding a vector database (e.g., Pinecone at \$70/month or a self-hosted pgvector instance requiring a larger cloud VM at \$25--50/month) would further increase costs.

The Featherless.ai free tier provides approximately 200 requests/day, which is sufficient for the expected query volume. If usage exceeds this, upgrading to a paid tier or switching to a self-hosted Llama 3.1 instance on a \$20/month cloud GPU would still keep costs under \$25/month.

%=====================================================================
\chapter{What Was Deliberately Not Built}
%=====================================================================

\section{No Vector Embeddings (pgvector)}

As discussed in Chapter 3, vector embeddings were deliberately excluded. The two-stage pipeline (regex keyword filter + LLM ranking) achieves the same goal---finding relevant contacts---with less complexity, lower cost, and better explainability. The 30--50 candidate limit from Stage 1 fits comfortably within the LLM's context window, making vector similarity search unnecessary.

\section{No LangChain}

LangChain is a popular framework for building LLM applications, providing abstractions for prompts, chains, agents, and tools. It was deliberately not used for three reasons:

\begin{enumerate}[itemsep=3pt]
    \item \textbf{Over-abstraction.} The matching pipeline has exactly one LLM call with a single prompt. LangChain's chain and agent abstractions add complexity without benefit for this use case.
    \item \textbf{Dependency weight.} LangChain adds approximately 50+ transitive dependencies. The current system uses only \texttt{httpx} for HTTP calls to Featherless.ai.
    \item \textbf{Debugging transparency.} Raw HTTP calls with explicit prompt construction make it easy to inspect exactly what is sent to the LLM and what is received. LangChain's abstraction layers obscure this.
\end{enumerate}

The agent module (\texttt{agent.py}) is 130 lines of Python that directly constructs a prompt, sends an HTTP POST to Featherless.ai, parses the JSON response, and handles errors. This is simpler, more maintainable, and more debuggable than an equivalent LangChain implementation.

\section{No Authentication}

The API currently has no authentication. This is acceptable for Phase 1 and 2 (local development and testing) but must be addressed before production deployment. Phase 3 will add API key authentication or JWT-based auth for the webhook endpoint.

\section{No WhatsApp Integration}

The system currently accepts queries via HTTP POST to \texttt{/match} and returns formatted WhatsApp-ready text. The actual WhatsApp message flow---receiving a message via the WhatsApp Business API, processing it through the matching engine, and sending the response back---is planned for Phase 3. The \texttt{messages} table in the database schema is already prepared for this.

%=====================================================================
\chapter{Challenges and How They Were Solved}
%=====================================================================

\section{Docker Desktop Not Running After System Restart}

\textbf{Problem:} After the development machine restarted, Docker Desktop was not automatically started. This caused the PostgreSQL container to be unavailable, which prevented the FastAPI server from starting (connection refused on port 5432).

\textbf{Solution:} Docker Desktop must be started manually before running the application. This is a known limitation of the development environment and will be addressed in Phase 4 by deploying to Oracle Cloud, where the database will run as a managed service or always-on container.

\section{Streamlit Deprecation Warnings}

\textbf{Problem:} Streamlit 1.42+ deprecated the \texttt{use\_container\_width} parameter in favour of \texttt{width} (with values \texttt{`stretch'} or \texttt{`content'}). The dashboard initially used \texttt{use\_container\_width=True} throughout, generating deprecation warnings on every page load.

\textbf{Solution:} All instances of \texttt{use\_container\_width=True} were replaced with \texttt{width=`stretch'} in \texttt{plotly\_chart} and \texttt{dataframe} calls. This eliminated the warnings and future-proofed the code.

\section{LLM Rate Limiting}

\textbf{Problem:} The Featherless.ai free tier imposes rate limits. When exceeded, the API returns HTTP 429 (Too Many Requests).

\textbf{Solution:} The agent implements retry logic with exponential backoff: 5 seconds, 10 seconds, then 15 seconds across 3 attempts. If all retries fail, a fallback response is returned with a user-friendly message. This ensures the API never returns a 500 error to the client.

\section{LLM JSON Parsing Failures}

\textbf{Problem:} Occasionally, the LLM returns malformed JSON (e.g., trailing commas, unescaped characters, or extra text outside the JSON block).

\textbf{Solution:} The agent wraps \texttt{json.loads()} in a try/except block. If parsing fails, the error is caught by the outer exception handler, which triggers the retry/fallback logic. In practice, with temperature 0.1 and a well-structured prompt, JSON parsing failures are rare (approximately 1 in 50 requests).

\section{Server Processes Terminating}

\textbf{Problem:} Both the FastAPI server and Streamlit dashboard run as foreground processes in terminal windows. When the terminal sessions end or the machine restarts, both servers stop.

\textbf{Solution:} For development, restart commands are documented. For production (Phase 4), both services will run as systemd services or Docker containers with restart policies (\texttt{restart: always}), ensuring they survive machine restarts.

%=====================================================================
\chapter{Phase 3 and 4 Roadmap}
%=====================================================================

\section{Phase 3: WhatsApp Integration}

Phase 3 will connect the matching engine to WhatsApp, enabling the full end-to-end flow:

\begin{enumerate}[itemsep=3pt]
    \item \textbf{WhatsApp Business API setup:} Register a business profile, configure a webhook endpoint, and obtain the necessary API credentials.
    \item \textbf{Webhook receiver:} A new FastAPI endpoint (\texttt{POST /webhook}) will receive incoming WhatsApp messages, validate the signature, extract the message text and sender number, and store the message in the \texttt{messages} table.
    \item \textbf{Message processing:} The webhook handler will call \texttt{MatchingEngine.match()} with the message text and receive the formatted response.
    \item \textbf{Response delivery:} The formatted response will be sent back to the user via the WhatsApp Business API.
    \item \textbf{Match history:} Each match result will be recorded in the \texttt{match\_history} table with the message ID, contact IDs, rank positions, confidence scores, and reasoning.
\end{enumerate}

\section{Phase 4: Production Deployment}

Phase 4 will harden the system for production use:

\begin{enumerate}[itemsep=3pt]
    \item \textbf{SSL/TLS:} Configure HTTPS for the API endpoint using Let's Encrypt or a purchased certificate.
    \item \textbf{Oracle Cloud deployment:} Deploy the FastAPI application and PostgreSQL database to Oracle Cloud Free Tier (4 ARM cores, 24 GB RAM, 200 GB storage).
    \item \textbf{Monitoring:} Add health check endpoints, request logging, and error alerting.
    \item \textbf{Load testing:} Verify the system can handle expected query volumes (100--200/day) with acceptable response times.
    \item \textbf{Authentication:} Add API key or JWT-based authentication to protect the endpoints.
\end{enumerate}

\section{Future Enhancements}

\begin{itemize}[itemsep=3pt]
    \item \textbf{Feedback loop:} Allow users to rate matches (accepted/rejected). Use this feedback to fine-tune the LLM prompt or adjust confidence thresholds over time.
    \item \textbf{Contact scoring:} Track which contacts are most frequently matched and accepted. Use this to dynamically adjust ranking weights.
    \item \textbf{Admin dashboard:} A protected web interface for managing contacts, viewing match history, and monitoring system health.
    \item \textbf{Multi-turn conversations:} Allow the LLM to ask clarifying questions and refine results based on user responses, rather than returning a single clarification question and ending the interaction.
\end{itemize}

%=====================================================================
\chapter{Conclusion}
%=====================================================================

The AI Middleman project demonstrates that a practical, production-quality AI matching system can be built at near-zero cost using open-weight models and thoughtful architecture. The two-stage pipeline---keyword filter for scale, LLM for intelligence---solves the fundamental constraint of context window limits while delivering ranked, explained matches in under 3 seconds.

The architectural decisions were driven by pragmatism: PostgreSQL for reliable relational storage, asyncpg for non-blocking database access, Featherless.ai's free tier for zero-cost LLM access, and a deliberate avoidance of vector embeddings and LangChain in favour of simpler, more transparent alternatives. Each decision was made by asking ``what is the simplest thing that works?'' rather than ``what is the most sophisticated approach available?''

The system currently handles 50,000 contacts, returns ranked matches with confidence scores and human-readable reasoning, and formats output as WhatsApp-ready messages. The Streamlit dashboard provides visualisation and manual testing capabilities. The test suite covers the critical pipeline components.

Phase 3 will add WhatsApp integration, completing the end-to-end flow from incoming message to outgoing match response.

\end{document}
